% Generated draft: requires mathematical and source review.
\documentclass[11pt]{article}
\usepackage[margin=25mm]{geometry}
\usepackage{fontspec}
\setmainfont{Latin Modern Roman}
\usepackage{amsmath,amssymb,mathtools,unicode-math}
\setmathfont{Latin Modern Math}
\usepackage{enumitem,needspace,xurl,hyperref}
\setlength{\parindent}{0pt}
\setlength{\parskip}{5pt}
\setlength{\emergencystretch}{3em}
\newcommand{\N}{\mathbb{N}}
\newcommand{\Z}{\mathbb{Z}}
\newcommand{\Q}{\mathbb{Q}}
\newcommand{\R}{\mathbb{R}}
\newcommand{\C}{\mathbb{C}}
\newcommand{\F}{\mathbb{F}}
\newcommand{\A}{\mathbb{A}}
\newcommand{\PP}{\mathbb{P}}
\newcommand{\E}{\mathbb{E}}
\newcommand{\eps}{\varepsilon}
\newcommand{\dd}{\,\mathrm{d}}
\begin{document}
\section*{UnsolvedMath ID: 30006037}
\textbf{Source problem code:} catalogue entry
\subsection{Definitions and mathematical statement}
Let \(n\) be a parameter for which the constructions in the cited report are defined. The supplied record does not specify its admissible values or the ground ring. Write \(C=\mathrm{EGraphs}_n^c\) for the report's chain cooperad and \(\mathfrak g=\mathrm{EGC}_{2n}\) for its \(L_\infty\)-algebra acting on \(C\). Their graph constructions, differentials, gradings, and action are not defined in the supplied excerpt and must be taken from the source before this formulation can be made self-contained.

Write \(\operatorname{Def}(C,C)\) for the deformation complex of the object \(C\) in the category of chain cooperads, denoted \(\operatorname{Dfm}(C,C)\) in the original excerpt. Its precise deformation-theoretic construction is not supplied. The notation \(E_k\) refers to the source's higher-algebra structure; the chain-operad model is unspecified. The source describes \(U_{E_{n+2}}\) as the enveloping \(E_{n+2}\)-algebra functor from \(L_\infty\)-algebras with divided powers to \(E_{n+2}\)-algebras. The divided-power qualification is part of the proposed statement.

\textbf{Conjecture, in the intended setting of the source.} The deformation complex \(\operatorname{Def}(C,C)\) admits a natural \(E_{n+2}\)-algebra structure, and the specified action of \(\mathfrak g\) on \(C\) induces a quasi-isomorphism
\[
U_{E_{n+2}}(\mathrm{EGC}_{2n})
\xrightarrow{\ \sim\ }
\operatorname{Def}(\mathrm{EGraphs}_n^c,\mathrm{EGraphs}_n^c).
\]
Here a quasi-isomorphism means a map inducing isomorphisms on the homology of the underlying complexes. The proposed map is to be compatible with the asserted higher-algebra structures. However, the supplied excerpt does not specify the category in which this compatibility is expressed, the meaning of naturality, the necessary shifts or completions, or the hypotheses on \(n\) and the ground ring. This draft therefore records the source's conjectural assertion with these qualifications unresolved; it does not assert the conjecture for an invented parameter range or coefficient system.
\subsection{Short English statement}
The conjecture asks whether the graph-cooperad deformation complex has a natural \(E_{n+2}\)-algebra structure and whether the specified graph \(L_\infty\)-action yields the displayed enveloping-algebra quasi-isomorphism. The source explicitly includes divided powers. The parameter range, coefficients, graph constructions, and precise compatibility remain unspecified in the supplied excerpt.
\subsection{Sources}
\begin{itemize}
\item Local UnsolvedMath registry, \texttt{data/unsolvedmath/problems.json}, supplied complete record with permanent ID \(30006037\), exact source code \texttt{catalogue entry}, title \emph{Higher Algebra Structures on Graph-Cooperad Deformations}. The record assigns proposed year 2024 and supplies no individual proposer.
\item Supplied primary citation: \emph{Homotopical Algebra and Higher Structures} (2025), Oberwolfach Reports, DOI 10.4171/owr/2024/39, pp.~2251--2308. \url{https://doi.org/10.4171/owr/2024/39}. The supplied literature-source label calls this \emph{Oberwolfach Report 39 (2024)}; the differing year labels are retained without reconciliation.
\item Supplied literature source: Fresse, \emph{Homotopical Algebra and Higher Structures}. \url{https://pro.univ-lille.fr/fileadmin/user_upload/pages_pros/benoit_fresse/Manuscripts/owr-BFresse-2024.pdf}. This reference label does not establish individual attribution of the conjecture.
\item Supplied literature assessment dated \textbf{2026-08-22}: the natural \(E_{n+2}\)-structure and enveloping-algebra quasi-isomorphism are described as explicitly conjectural in the primary report, with no later proof identified by that assessment. This is supplied dated curation derived from the Oberwolfach material, not independently verified literature status. The registry labels the problem open and the statement \texttt{corrected\_verified}, and reports that the official report was checked and its notation repaired; those are supplied editorial claims, not verification performed for this draft.
\end{itemize}
\end{document}
